% META: {"id":"TCOMPL-MF-EX-002","chapitre":"TCOMPL-MODELES-FONCTION","type_objet":"exercice","capacites":["TCOMPL-MF-C2"],"capacites_codes":["C2"],"methodes":[],"parcours":2,"competences":["raisonner"],"duree_min":15,"mode_creation":"ex_nihilo","sources_inspiration":[],"fichier_tex":"chapitres/TCOMPL-MODELES-FONCTION/exercices/TCOMPL-MF-EX-002.tex","corrige_tex":"chapitres/TCOMPL-MODELES-FONCTION/corriges/TCOMPL-MF-CO-002.tex","status":"approved","origin":"generated","status_history":["generated","audited","accepted","approved"],"release_acceptance":"RELEASE_OWNER_BATCH_ACCEPTANCE","acceptance_closure_digest":"sha256:9b3ccf9a81c5520fbb7e03b7b2d3e7bf2057a02834b6d908bf3be5f49f3b553f","acceptance_receipt_digest":"sha256:4fad86f0282e0fa4e0419cca1ad6379ae7af5a280c04a4a90880f90a1c044242"}
% BEGIN-VERIFY
% from sympy import symbols, solve, Eq
% x = symbols('x', positive=True)
% f = x + 16/x
% assert len(solve(Eq(f, 10), x)) == 2
% assert solve(Eq(f, 8), x) == [4]
% assert solve(Eq(f, 7), x) == []
% END-VERIFY
\begin{exercice}{TCOMPL-MF-EX-002}{2}{15}
On reprend $f(x) = x + \dfrac{16}{x}$ (décroissante sur $]0,4]$ de $+\infty$ à $8$, puis croissante sur $[4,+\infty[$ de $8$ à $+\infty$). Sans chercher à résoudre explicitement les équations, déterminer, en justifiant à l'aide du tableau de variation, le nombre de solutions de :
\begin{enumerate}
  \item $f(x) = 10$ ;
  \item $f(x) = 8$ ;
  \item $f(x) = 7$.
\end{enumerate}
\end{exercice}
